\section{Results, impact, and closing}
\label{sec:results}

\subsection{Evaluation protocol (slide 22)}

\begin{whybox}
Defines what was measured before you flash numbers. Honesty about scope builds trust.
\end{whybox}

\begin{itemize}
  \item Stratified \(10\%\) holdout from Triagegeist: \(N_{\mathrm{eval}}=8{,}000\) (seed \(42\)).
  \item Same English complaint texts and nurse acuity labels as the training distribution.
  \item Metrics: accuracy, macro-F1, per-class precision/recall/F1, latency (p50, p95, mean).
\end{itemize}

\begin{remark}[What the numbers mean---and do not mean]
Holdout scores measure agreement with historical nurse labels on \textbf{in-distribution} Triagegeist text.
They are \textbf{not} prospective waiting-time or outcome trials at \examplehospital{} \citep{stewart2023nlp}.
Say this out loud if asked ``does this prove patients get better care at KNUST?''
\end{remark}

\begin{saybox}
``We evaluate on eight thousand held-out Triagegeist complaints. High scores mean the model matches nurse labels on that benchmark---not that we have finished a clinical trial at KNUST.''
\end{saybox}

\subsection{Holdout metrics and SMOTE trade-off (slide 23)}

\begin{whybox}
Shows headline performance and the safety trade-off for the rare Red class.
\end{whybox}

\begin{center}
\small
\begin{tabular}{@{}lrr@{}}
\toprule
Metric & Baseline & SMOTE \\
\midrule
Accuracy & \(0.9992\) & \(0.9985\) \\
Macro-F1 & \(0.9977\) & \(0.9954\) \\
Recall Level~1 (Red) & \(0.9814\) & \(1.0000\) \\
Precision Level~1 (Red) & \(1.0000\) & \(0.9641\) \\
Latency p50 (ms) & \(100.30\) & \(95.50\) \\
\bottomrule
\end{tabular}
\end{center}

\paragraph{How to explain the trade-off.}
\begin{itemize}
  \item \textbf{Baseline (default):} nearly perfect overall accuracy; perfect Red \emph{precision}
    (when it says Red, it was Red) but Red \emph{recall} \(0.9814\) (a few Reds missed).
  \item \textbf{SMOTE} \citep{chawla2002smote}: oversamples minority classes in training.
    Red recall becomes \(100\%\) (safer net for critical cases) with a small drop in Red precision
    and a tiny drop in overall accuracy/F1.
\end{itemize}

\paragraph{Recall vs precision in one sentence.}
Recall (Red) = among true Reds, how many did we catch?
Precision (Red) = among predicted Reds, how many were truly Red?
For critical triage, missing a Red (low recall) is often worse than an extra Orange/Yellow alert.

\begin{saybox}
``Baseline is extremely accurate with perfect Red precision. SMOTE pushes Red recall to one hundred percent for a safer clinical net, at a small precision cost.''
\end{saybox}

\subsection{Headline chart and confusion matrix (slides 24--25)}

\begin{whybox}
Visual confirmation of the table; confusion matrix shows where errors (if any) sit.
\end{whybox}

\paragraph{Bar chart.}
Same story as the table: baseline vs SMOTE on accuracy, macro-F1, Level~1 recall.
Point at Red recall rising under SMOTE.

\paragraph{Confusion matrix (baseline).}
Rows = true labels, columns = predicted labels (or as labelled on the figure).
A strong diagonal means almost all complaints are classified correctly on the holdout.
If asked about off-diagonal cells: those are the rare mistakes; relate them to recall/precision.

\begin{saybox}
``The confusion matrix is nearly diagonal on the holdout---errors are rare. That matches the \(0.999\) accuracy.''
\end{saybox}

\subsection{Latency and nurse vs chatbot (slide 26)}

\begin{whybox}
Answers ``is this fast enough for intake?'' and contrasts queue time with model time.
\end{whybox}

\begin{itemize}
  \item Median inference \(\approx 100\)\,ms per complaint on evaluation hardware.
  \item Nurse path: minutes to hours to first colour (queue).
  \item Chatbot: English text \(\to\) colour + pathway card in about a tenth of a second.
  \item Nurse still performs TEWS and can override \citep{Mitchell2023}.
\end{itemize}

\begin{saybox}
``The model responds in about one hundred milliseconds. The nurse still owns vitals and the final clinical decision.''
\end{saybox}

\subsection{Anticipated intake workflow (slide 27)}

\begin{whybox}
Closes the story visually: traditional long queue vs chatbot-assisted path with TEWS in parallel.
\end{whybox}

\textbf{Traditional:} arrive \(\to\) queue \(\to\) nurse colour \(\to\) often unclear destination (minutes to hours).

\textbf{Chatbot-assisted:} complaint \(\to\) classify \(\to\) pathway card (\(\approx 100\)\,ms), nurse TEWS in parallel.

\begin{saybox}
``We are not removing the nurse. We insert a fast text-based colour and pathway while vitals continue as usual.''
\end{saybox}

\subsection{Anticipated impact (slide 28)}

\begin{whybox}
Four plain claims for a mixed audience. No jargon needed here.
\end{whybox}

\begin{enumerate}
  \item \textbf{Faster start:} suggest urgency as soon as the problem is described.
  \item \textbf{Safer for critical patients:} designed to catch the most urgent cases (tie to high Red recall).
  \item \textbf{Clearer next steps:} pathway card says where to go and what to do.
  \item \textbf{Nurses stay in charge:} assistive tool; vitals and override remain with staff.
\end{enumerate}

\begin{saybox}
``Faster first signal, fewer missed critical cases on the benchmark, clearer routing, and nurses remain in control.''
\end{saybox}

\subsection{Selected references (slide 29)}

\begin{whybox}
You will not read this aloud. Know the big names if asked: Lee (BioBERT), Stewart (NLP triage review),
SATS manual, Chawla (SMOTE), Triagegeist, plus local/context citations.
\end{whybox}

\subsection{One-minute closing script}

\begin{saybox}[Closing]
``We built an examinable pipeline from English chief complaint text to SATS colour and a KNUST pathway card,
with BioBERT mathematics, a clinical gate, and strong holdout agreement on Triagegeist.
The nurse keeps TEWS and override. That is the contribution of this project.''
\end{saybox}
